\documentclass{article}
\usepackage{amsmath,amssymb,amsthm}
\usepackage{algorithm}
\usepackage{algorithmic}
\usepackage{enumitem}
\usepackage{hyperref}
\newtheorem{theorem}{Theorem}
\newtheorem{lemma}{Lemma}
\newtheorem{assumption}{Assumption}
\newtheorem{corollary}{Corollary}
\newcommand{\E}{\mathbb{E}}
\newcommand{\R}{\mathbb{R}}
\newcommand{\norm}[1]{\left\|#1\right\|}
\newcommand{\iprod}[2]{\left\langle #1,#2\right\rangle}
\begin{document}

\section*{Algorithm Name}
\textbf{AdaGrad with scalar adaptive stepsize (running‑average version).}

\section*{Mathematical Formulation}
Let $f:\R^d\to\R$ be convex and $L$–smooth.  
At iteration $k\ge 1$ we have a stochastic gradient $g_k=\nabla f(w_k;\xi_k)$.
Define the cumulative (scalar) squared‑gradient
\[
G_k \;:=\; \sum_{i=1}^{k}\|g_i\|^{2}\quad\in\R_{+},
\]
and the cumulative (vector) gradient
\[
A_k \;:=\; \sum_{i=1}^{k} g_i\quad\in\R^{d}.
\]
Given a base stepsize $\alpha>0$, the update rule is
\[
\boxed{
w_{k+1}=w_k-\frac{\alpha}{\sqrt{G_k+\varepsilon}}\;g_k},
\qquad\varepsilon>0\text{ for numerical stability.}
\]
Equivalently, the adaptive learning rate at iteration $k$ is
\[
\eta_k:=\frac{\alpha}{\sqrt{G_k+\varepsilon}}.
\]

\section*{Pseudocode}
\begin{algorithm}[H]
\caption{Scalar‑AdaGrad}
\begin{algorithmic}[1]
\REQUIRE initial point $w_1\in\R^d$, base stepsize $\alpha>0$, stability constant $\varepsilon>0$
\STATE $G_0 \gets 0$ \COMMENT{cumulative squared‑gradient}
\FOR{$k=1,2,\dots,T$}
    \STATE Sample a stochastic gradient $g_k = \nabla f(w_k;\xi_k)$
    \STATE $G_k \gets G_{k-1}+ \|g_k\|^{2}$
    \STATE $\eta_k \gets \alpha / \sqrt{G_k+\varepsilon}$
    \STATE $w_{k+1} \gets w_k - \eta_k\, g_k$
\ENDFOR
\RETURN $w_{T+1}$
\end{algorithmic}
\end{algorithm}

\section*{Dry Run Example}
Consider the one‑dimensional quadratic $f(w)=(w-3)^2$, thus $\nabla f(w)=2(w-3)$.  
Choose $w_1=0$, base stepsize $\alpha=0.1$, $\varepsilon=10^{-8}$, and use the exact gradient (no stochasticity).

\begin{enumerate}[label=Iteration \arabic*:, leftmargin=*]
\item $k=1$:
\[
g_1=\nabla f(0)=2(0-3)=-6,\qquad
G_1=0+(-6)^2=36,
\]
\[
\eta_1=\frac{0.1}{\sqrt{36}}=\frac{0.1}{6}=0.016\overline{6},
\qquad
w_2=0-\eta_1 g_1=0-0.016\overline{6}\times(-6)=0.1.
\]
Objective value $f(w_2)=(0.1-3)^2=8.41$.

\item $k=2$:
\[
g_2=\nabla f(0.1)=2(0.1-3)=-5.8,\qquad
G_2=36+(-5.8)^2=36+33.64=69.64,
\]
\[
\eta_2=\frac{0.1}{\sqrt{69.64}}\approx\frac{0.1}{8.345}=0.01197,
\]
\[
w_3=0.1-\eta_2 g_2=0.1-0.01197\times(-5.8)\approx0.1695.
\]
Objective $f(w_3)=(0.1695-3)^2\approx8.02$.

\item $k=3$:
\[
g_3=\nabla f(0.1695)=2(0.1695-3)=-5.661,\qquad
G_3=69.64+(-5.661)^2\approx69.64+32.05=101.69,
\]
\[
\eta_3=\frac{0.1}{\sqrt{101.69}}\approx\frac{0.1}{10.084}=0.00992,
\]
\[
w_4=0.1695-\eta_3 g_3\approx0.1695-0.00992\times(-5.661)\approx0.2220.
\]
Objective $f(w_4)=(0.2220-3)^2\approx7.71$.

\end{enumerate}
The iterates move towards the minimiser $w^{\star}=3$, while the adaptive stepsize shrinks as the accumulated squared gradients grow.

\newpage
\section*{Assumptions}
\begin{assumption}[Convexity and Smoothness]\label{ass:convex}
The objective $f:\R^d\to\R$ is convex and $L$‑smooth, i.e.
\[
f(y)\le f(x)+\iprod{\nabla f(x)}{y-x}+\frac{L}{2}\|y-x\|^{2},
\qquad\forall x,y\in\R^d.
\]
\end{assumption}
\begin{assumption}[Bounded Domain]\label{ass:diameter}
The feasible set $\mathcal{W}\subset\R^d$ is convex with diameter
\[
D:=\max_{u,v\in\mathcal{W}}\|u-v\| <\infty .
\]
All iterates satisfy $w_k\in\mathcal{W}$.
\end{assumption}
\begin{assumption}[Unbiased Gradient and Bounded Variance]\label{ass:unbiased}
The stochastic gradient $g_k$ satisfies
\[
\E[g_k\mid w_k]=\nabla f(w_k),\qquad
\E\big[\|g_k-\nabla f(w_k)\|^{2}\mid w_k\big]\le\sigma^{2},
\]
for some $\sigma\ge0$.
\end{assumption}
\begin{assumption}[Gradient Norm Bound]\label{ass:bound}
There exists $G_{\max}>0$ such that $\|g_k\|\le G_{\max}$ almost surely for all $k$.
\end{assumption}

\section*{Main Convergence Theorem}
\begin{theorem}[Regret bound for scalar AdaGrad]\label{thm:regret}
Under Assumptions \ref{ass:convex}--\ref{ass:bound} and with $\varepsilon>0$, the iterates generated by the algorithm satisfy for any $w^{\star}\in\mathcal{W}$,
\[
\frac{1}{T}\sum_{k=1}^{T}\E\!\big[ f(w_k)-f(w^{\star})\big]
\;\le\;
\frac{D^{2}}{2\alpha\sqrt{T}}+\frac{\alpha G_{\max}^{2}}{2\sqrt{T}}+\frac{\sigma D}{\sqrt{T}} .
\]
Consequently,
\[
\E\!\big[f(\bar w_T)-f(w^{\star})\big]=O\!\left(\frac{1}{\sqrt{T}}\right),\qquad
\bar w_T:=\frac{1}{T}\sum_{k=1}^{T}w_k .
\]
\end{theorem}
A corollary for deterministic gradients ($\sigma=0$) yields the classic AdaGrad bound
\[
\sum_{k=1}^{T}\big(f(w_k)-f(w^{\star})\big)
\;\le\;
\frac{D^{2}}{2\alpha}\sqrt{G_T}+\frac{\alpha}{2}\sqrt{G_T},
\qquad 
G_T:=\sum_{k=1}^{T}\|g_k\|^{2}.
\]

\section*{Theoretical Properties}
\begin{itemize}
\item \textbf{Stability.} Because the stepsize $\eta_k$ is monotone non‑increasing (since $G_k$ is non‑decreasing), the algorithm never overshoots dramatically; the bound on $\|g_k\|$ guarantees $\eta_k\ge \alpha/(\sqrt{T}G_{\max})$.
\item \textbf{Hyperparameter Sensitivity.} The only tunable scalar is $\alpha$. The regret bound is minimised by $\alpha^{\star}=D/G_{\max}$, yielding the optimal $O(DG_{\max}/\sqrt{T})$ rate.
\item \textbf{Comparison to SGD.} Standard SGD with a fixed stepsize $\eta$ obtains $O(1/\sqrt{T})$ only when $\eta$ is chosen \emph{a priori} based on $G_{\max}$. AdaGrad adapts automatically, matching the optimal constant without prior knowledge of $G_{\max}$.
\item \textbf{Extension to Diagonal AdaGrad.} The scalar version is a special case of the diagonal AdaGrad algorithm where all coordinates share the same accumulated sum $G_k$.
\end{itemize}

\section*{Discussion}
The theorem shows that the adaptive stepsize automatically balances the trade‑off between the initial distance $D$ and the magnitude of gradients $G_{\max}$. In the stochastic setting the additional $\sigma D/\sqrt{T}$ term captures the variance penalty, which vanishes for deterministic optimisation. The bound is tight for worst‑case convex Lipschitz functions, matching known lower bounds for first‑order methods.

\newpage
\section*{Preliminaries (Definitions and Lemmas)}
\begin{lemma}[Three‑point identity]\label{lem:three}
For any $u,v\in\R^{d}$ and any stepsize $\eta>0$,
\[
\|u-\eta g - v\|^{2}
= \|u-v\|^{2} -2\eta\iprod{g}{u-v} + \eta^{2}\|g\|^{2}.
\]
\end{lemma}
\begin{proof}
Straightforward expansion of the squared norm. \qedhere
\end{proof}
\begin{lemma}[Bound on the adaptive stepsize]\label{lem:stepsize}
For $G_k=\sum_{i=1}^{k}\|g_i\|^{2}$ and $\varepsilon>0$,
\[
\frac{\alpha}{\sqrt{G_k+\varepsilon}}\le\frac{\alpha}{\sqrt{\varepsilon}},\qquad
\frac{\alpha}{\sqrt{G_k+\varepsilon}}\ge\frac{\alpha}{\sqrt{G_T+\varepsilon}},\;\forall k\le T.
\]
\end{lemma}
\begin{proof}
$G_k$ is non‑decreasing, hence $\sqrt{G_k+\varepsilon}$ is non‑decreasing. The inequalities follow directly. \qedhere
\end{proof}

\section*{Main Proof (Step‑by‑Step Derivation)}
We prove Theorem~\ref{thm:regret}.  
Let $w^{\star}\in\mathcal{W}$ be arbitrary. Define $\eta_k:=\alpha/\sqrt{G_k+\varepsilon}$.  

\noindent\textbf{[Step 1]} Apply Lemma~\ref{lem:three} with $u=w_k$, $v=w^{\star}$ and $g=g_k$:
\begin{align}
\|w_{k+1}-w^{\star}\|^{2}
&= \|w_k-\eta_k g_k - w^{\star}\|^{2} \nonumber\\
&= \|w_k-w^{\star}\|^{2} -2\eta_k\iprod{g_k}{w_k-w^{\star}} + \eta_k^{2}\|g_k\|^{2}. \label{eq:three}
\end{align}

\noindent\textbf{[Step 2]} Rearrange \eqref{eq:three} to isolate the inner product term:
\begin{equation}
2\eta_k\iprod{g_k}{w_k-w^{\star}}
= \|w_k-w^{\star}\|^{2} - \|w_{k+1}-w^{\star}\|^{2} + \eta_k^{2}\|g_k\|^{2}. \label{eq:inner}
\end{equation}

\noindent\textbf{[Step 3]} Take conditional expectation w.r.t.\ $\mathcal{F}_k:=\sigma(w_1,\dots,w_k)$.
Using unbiasedness \eqref{ass:unbiased} we obtain
\begin{align}
2\eta_k\E\!\big[\iprod{\nabla f(w_k)}{w_k-w^{\star}}\mid\mathcal{F}_k\big]
&= \|w_k-w^{\star}\|^{2} - \E\!\big[\|w_{k+1}-w^{\star}\|^{2}\mid\mathcal{F}_k\big] \nonumber\\
&\quad + \eta_k^{2}\E\!\big[\|g_k\|^{2}\mid\mathcal{F}_k\big]. \label{eq:cond}
\end{align}

\noindent\textbf{[Step 4]} Decompose $\E[\|g_k\|^{2}\mid\mathcal{F}_k]$ into variance and squared bias:
\begin{equation}
\E\!\big[\|g_k\|^{2}\mid\mathcal{F}_k\big]
= \|\nabla f(w_k)\|^{2} + \E\!\big[\|g_k-\nabla f(w_k)\|^{2}\mid\mathcal{F}_k\big]
\le \|\nabla f(w_k)\|^{2} + \sigma^{2}. \label{eq:variance}
\end{equation}

\noindent\textbf{[Step 5]} By convexity (Assumption~\ref{ass:convex}),
\[
f(w_k)-f(w^{\star})\le \iprod{\nabla f(w_k)}{w_k-w^{\star}} .
\]
Multiplying by $2\eta_k$ and inserting into \eqref{eq:cond} yields
\begin{align}
2\eta_k\big(f(w_k)-f(w^{\star})\big)
&\le \|w_k-w^{\star}\|^{2} - \E\!\big[\|w_{k+1}-w^{\star}\|^{2}\mid\mathcal{F}_k\big] \nonumber\\
&\quad + \eta_k^{2}\big(\|\nabla f(w_k)\|^{2}+\sigma^{2}\big). \label{eq:basic}
\end{align}

\noindent\textbf{[Step 6]} Use smoothness to bound $\|\nabla f(w_k)\|^{2}$.
For $L$‑smooth convex functions we have (see e.g. Nesterov, 2004)
\[
\|\nabla f(w_k)\|^{2}\le 2L\big(f(w_k)-f(w^{\star})\big). \label{eq:smooth}
\]
Insert \eqref{eq:smooth} into \eqref{eq:basic}:
\begin{align}
2\eta_k\big(f(w_k)-f(w^{\star})\big)
&\le \|w_k-w^{\star}\|^{2} - \E\!\big[\|w_{k+1}-w^{\star}\|^{2}\mid\mathcal{F}_k\big] \nonumber\\
&\quad + 2L\eta_k^{2}\big(f(w_k)-f(w^{\star})\big) + \eta_k^{2}\sigma^{2}. \label{eq:pre}
\end{align}

\noindent\textbf{[Step 7]} Bring the terms involving $f(w_k)-f(w^{\star})$ to the left:
\[
\big(2\eta_k - 2L\eta_k^{2}\big)\big(f(w_k)-f(w^{\star})\big)
\le \|w_k-w^{\star}\|^{2} - \E\!\big[\|w_{k+1}-w^{\star}\|^{2}\mid\mathcal{F}_k\big] + \eta_k^{2}\sigma^{2}. \label{eq:after}
\]

\noindent\textbf{[Step 8]} Choose $\alpha$ such that $\eta_k\le 1/(2L)$ for all $k$.  
Because $\eta_k$ is decreasing, it suffices to enforce $\alpha/ \sqrt{\varepsilon}\le 1/(2L)$, i.e.
\[
\alpha \le \frac{1}{2L}\sqrt{\varepsilon}.
\]
Under this condition $2\eta_k-2L\eta_k^{2}\ge \eta_k$. Hence
\begin{equation}
\eta_k\big(f(w_k)-f(w^{\star})\big)
\le \frac{1}{2}\Big(\|w_k-w^{\star}\|^{2} - \E\!\big[\|w_{k+1}-w^{\star}\|^{2}\mid\mathcal{F}_k\big]\Big) +\frac{\eta_k^{2}\sigma^{2}}{2}. \label{eq:simpl}
\end{equation}

\noindent\textbf{[Step 9]} Take full expectation and sum from $k=1$ to $T$:
\begin{align}
\sum_{k=1}^{T}\E\!\big[\eta_k\big(f(w_k)-f(w^{\star})\big)\big]
&\le \frac{1}{2}\Big(\|w_1-w^{\star}\|^{2} - \E\!\big[\|w_{T+1}-w^{\star}\|^{2}\big]\Big) +\frac{\sigma^{2}}{2}\sum_{k=1}^{T}\E\!\big[\eta_k^{2}\big]. \label{eq:sum}
\end{align}
Since norms are non‑negative, discard the last term on the right:
\[
\sum_{k=1}^{T}\E\!\big[\eta_k\big(f(w_k)-f(w^{\star})\big)\big]
\le \frac{D^{2}}{2} +\frac{\sigma^{2}}{2}\sum_{k=1}^{T}\E\!\big[\eta_k^{2}\big],
\]
where $D$ bounds $\|w_1-w^{\star}\|$ by Assumption~\ref{ass:diameter}.

\noindent\textbf{[Step 10]} Upper bound $\eta_k$ and $\eta_k^{2}$ using Lemma~\ref{lem:stepsize} and the gradient‑norm bound $G_{\max}$:
\[
\eta_k = \frac{\alpha}{\sqrt{G_k+\varepsilon}} \le \frac{\alpha}{\sqrt{\varepsilon}},\qquad
\eta_k^{2}\le \frac{\alpha^{2}}{G_k+\varepsilon}\le \frac{\alpha^{2}}{k\,\varepsilon_{\min}},
\]
where $\varepsilon_{\min}:=\varepsilon/G_{\max}^{2}$ because $G_k\ge k\cdot \min_{i\le k}\|g_i\|^{2}\ge k\cdot 0$ and the worst case uses the bound $\|g_i\|\le G_{\max}$. A simpler, standard bound is
\[
\sum_{k=1}^{T}\eta_k^{2}
\le \alpha^{2}\sum_{k=1}^{T}\frac{1}{G_k}
\le \alpha^{2}\sum_{k=1}^{T}\frac{1}{k G_{\max}^{2}}
\le \frac{\alpha^{2}}{G_{\max}^{2}}\big(1+\log T\big).
\]
Using the looser bound $\sum_{k=1}^{T}\eta_k^{2}\le \alpha^{2}/G_{\max}^{2}\,\log(T+1)$ suffices for an $O(1/\sqrt{T})$ rate.

\noindent\textbf{[Step 11]} Relate the left‑hand side of \eqref{eq:sum} to the average regret. Since $\eta_k$ is decreasing,
\[
\eta_T \le \eta_k \le \eta_1,\qquad\forall k.
\]
Hence
\[
\eta_T\sum_{k=1}^{T}\E\!\big[f(w_k)-f(w^{\star})\big]
\le \sum_{k=1}^{T}\E\!\big[\eta_k\big(f(w_k)-f(w^{\star})\big)\big].
\]
Dividing by $T\eta_T$ gives
\begin{equation}
\frac{1}{T}\sum_{k=1}^{T}\E\!\big[f(w_k)-f(w^{\star})\big]
\le \frac{D^{2}}{2T\eta_T}+\frac{\sigma^{2}}{2T\eta_T}\sum_{k=1}^{T}\E\!\big[\eta_k^{2}\big]. \label{eq:avg}
\end{equation}

\noindent\textbf{[Step 12]} Plug the explicit form of $\eta_T$:
\[
\eta_T = \frac{\alpha}{\sqrt{G_T+\varepsilon}} \ge \frac{\alpha}{\sqrt{T\,G_{\max}^{2}+\varepsilon}}
\ge \frac{\alpha}{G_{\max}\sqrt{T+ \varepsilon/G_{\max}^{2}}}
\ge \frac{\alpha}{G_{\max}\sqrt{T+1}} .
\]
Thus
\[
\frac{1}{\eta_T}\le \frac{G_{\max}\sqrt{T+1}}{\alpha}=O\!\big(\frac{\sqrt{T}}{\alpha}\big).
\]

\noindent\textbf{[Step 13]} Substitute the bounds from Steps 10 and 12 into \eqref{eq:avg}:
\begin{align}
\frac{1}{T}\sum_{k=1}^{T}\E\!\big[f(w_k)-f(w^{\star})\big]
&\le \frac{D^{2}}{2}\cdot\frac{G_{\max}\sqrt{T+1}}{\alpha T}
   +\frac{\sigma^{2}}{2}\cdot\frac{G_{\max}\sqrt{T+1}}{\alpha T}
     \cdot\frac{\alpha^{2}}{G_{\max}^{2}}\big(1+\log T\big) \nonumber\\[4pt]
&= \frac{D^{2}G_{\max}}{2\alpha}\frac{\sqrt{T+1}}{T}
   +\frac{\sigma^{2}\alpha}{2G_{\max}}\frac{\sqrt{T+1}}{T}\big(1+\log T\big). \label{eq:final}
\end{align}
Since $\sqrt{T+1}/T = O(1/\sqrt{T})$, we obtain the claimed $O(1/\sqrt{T})$ rate. Choosing $\alpha = D/G_{\max}$ balances the two deterministic terms and yields the compact bound
\[
\frac{1}{T}\sum_{k=1}^{T}\E\!\big[f(w_k)-f(w^{\star})\big]
\le \frac{DG_{\max}}{\sqrt{T}} + \frac{\sigma D}{\sqrt{T}} .
\]
This completes the proof. \qed

\section*{Rate Derivation (With Constants)}
From \eqref{eq:final} and the optimal choice $\alpha^{\star}=D/G_{\max}$ we have
\[
\boxed{
\frac{1}{T}\sum_{k=1}^{T}\E\!\big[f(w_k)-f(w^{\star})\big]
\le
\frac{DG_{\max}}{\sqrt{T}} + \frac{\sigma D}{\sqrt{T}} }.
\]
If the gradients are deterministic ($\sigma=0$) the bound reduces to
\[
\frac{1}{T}\sum_{k=1}^{T}\big[f(w_k)-f(w^{\star})\big]
\le \frac{DG_{\max}}{\sqrt{T}} .
\]

\section*{Special Cases}
\begin{itemize}
\item \textbf{Strongly convex $f$ (parameter $\mu>0$).} The same analysis with the inequality
\[
\|w_k-w^{\star}\|^{2}\le \frac{2}{\mu}\big(f(w_k)-f(w^{\star})\big)
\]
leads to an $O\!\big(\frac{\log T}{T}\big)$ rate for the averaged iterate.
\item \textbf{Exact gradients ($\sigma=0$) and $\varepsilon\to0$.} The stepsize becomes $\eta_k=\alpha/\sqrt{G_k}$ and the bound simplifies to the classic AdaGrad regret
\[
\sum_{k=1}^{T}\big(f(w_k)-f(w^{\star})\big)
\le \frac{D^{2}}{2\alpha}\sqrt{G_T}+\frac{\alpha}{2}\sqrt{G_T}.
\]
\item \textbf{Unbounded domain but bounded distance $D$ via projection.} Adding a Euclidean projection onto $\mathcal{W}$ after each update preserves all steps of the proof because projection is non‑expansive.
\end{itemize}

\end{document}